\documentclass[11pt]{article}
\usepackage[T1]{fontenc}
\usepackage{lmodern}
\usepackage[letterpaper,margin=1in,headheight=28pt]{geometry}
\usepackage{amsmath,amssymb,amsthm}
\usepackage{mathtools}
\usepackage{xcolor}
\usepackage[skins,breakable]{tcolorbox}
\usepackage{enumitem}
\usepackage{fancyhdr}
\usepackage{lastpage}

% Course and student details for the header.
\newcommand{\course}{MATH 221: Linear Algebra}
\newcommand{\psetnumber}{4}
\newcommand{\student}{Robin Park}
\newcommand{\duedate}{Friday, 2 October 2026}

\pagestyle{fancy}
\fancyhf{}
\fancyhead[L]{\textbf{\course}\\Problem Set \psetnumber}
\fancyhead[R]{\student\\Due \duedate}
\fancyfoot[C]{\thepage\ / \pageref{LastPage}}
\setlength{\parindent}{0pt}
\setlength{\parskip}{0.5em}

\definecolor{probblue}{RGB}{30,80,150}
\definecolor{solgreen}{RGB}{235,245,238}

% Problems are numbered automatically; the optional argument gives a short title.
\newcounter{problem}
\newenvironment{problem}[1][]{%
  \refstepcounter{problem}\par\bigskip
  {\color{probblue}\large\bfseries Problem \theproblem\if\relax\detokenize{#1}\relax\else\ (#1)\fi}\par\nopagebreak
}{\par}

% Solutions sit in a light box that can break across pages.
\newenvironment{solution}{%
  \begin{tcolorbox}[enhanced,breakable,colback=solgreen,colframe=solgreen,
    borderline west={3pt}{0pt}{green!45!black},arc=0pt,left=8pt]
  \textbf{Solution.}\ }{\hfill$\blacksquare$\end{tcolorbox}}

\newcommand{\R}{\mathbb{R}}
\DeclarePairedDelimiter{\norm}{\lVert}{\rVert}

\begin{document}

\begin{problem}[Solving a system]
Find all solutions of the linear system
\[
\begin{aligned}
  x + 2y - z &= 3,\\
  2x + 5y + z &= 8,\\
  -x - y + 4z &= -1.
\end{aligned}
\]
\end{problem}

\begin{solution}
Row reduce the augmented matrix:
\[
\left[\begin{array}{ccc|c} 1 & 2 & -1 & 3\\ 2 & 5 & 1 & 8\\ -1 & -1 & 4 & -1 \end{array}\right]
\sim
\left[\begin{array}{ccc|c} 1 & 2 & -1 & 3\\ 0 & 1 & 3 & 2\\ 0 & 1 & 3 & 2 \end{array}\right]
\sim
\left[\begin{array}{ccc|c} 1 & 0 & -7 & -1\\ 0 & 1 & 3 & 2\\ 0 & 0 & 0 & 0 \end{array}\right].
\]
So $z = t$ is free, $y = 2 - 3t$ and $x = -1 + 7t$. The solution set is the line
\[
  \begin{pmatrix} x\\ y\\ z \end{pmatrix} = \begin{pmatrix} -1\\ 2\\ 0 \end{pmatrix} + t\begin{pmatrix} 7\\ -3\\ 1 \end{pmatrix}, \qquad t \in \R.
\]
\end{solution}

\begin{problem}[Eigenvalues]
Let $A = \begin{pmatrix} 4 & 1\\ 2 & 3 \end{pmatrix}$.
\begin{enumerate}[label=(\alph*)]
  \item Compute the eigenvalues of $A$.
  \item Find an eigenvector for each eigenvalue.
  \item Is $A$ diagonalisable? Justify your answer.
\end{enumerate}
\end{problem}

\begin{solution}
\begin{enumerate}[label=(\alph*)]
  \item $\det(A - \lambda I) = (4-\lambda)(3-\lambda) - 2 = \lambda^2 - 7\lambda + 10 = (\lambda-2)(\lambda-5)$, so $\lambda_1 = 2$ and $\lambda_2 = 5$.
  \item For $\lambda_1 = 2$: $(A - 2I)v = 0$ gives $v_1 = (1, -2)^{T}$. For $\lambda_2 = 5$: $v_2 = (1, 1)^{T}$.
  \item Yes. The two eigenvalues are distinct, so $v_1, v_2$ are linearly independent and $A = PDP^{-1}$ with $P = [v_1\ v_2]$ and $D = \operatorname{diag}(2, 5)$.
\end{enumerate}
\end{solution}

\begin{problem}
Prove that for all vectors $u, v \in \R^n$,
\[
  \norm{u + v}^2 + \norm{u - v}^2 = 2\norm{u}^2 + 2\norm{v}^2.
\]
\end{problem}

\begin{solution}
Expand both norms using the dot product:
\begin{align*}
  \norm{u+v}^2 &= u\cdot u + 2\,u\cdot v + v\cdot v,\\
  \norm{u-v}^2 &= u\cdot u - 2\,u\cdot v + v\cdot v.
\end{align*}
Adding the two lines, the cross terms cancel and we obtain $2\,u\cdot u + 2\,v\cdot v = 2\norm{u}^2 + 2\norm{v}^2$.
\end{solution}

\begin{problem}[Your turn]
% Copy this block for more problems.
State the problem here.
\end{problem}

\begin{solution}
Write your solution here.
\end{solution}

\end{document}
